\documentclass[10pt,a4paper]{article}

\usepackage[utf8]{inputenc}
\usepackage[T1]{fontenc}
\usepackage{lmodern}
% Portuguese hyphenation when babel-portuges is installed (it is on Overleaf);
% the accents below are plain TeX commands, so it also compiles without it.
\IfFileExists{portuges.ldf}{\usepackage[portuguese]{babel}}{\emergencystretch=3em}
\usepackage[a4paper,margin=2cm,top=2cm,bottom=2cm]{geometry}
\usepackage{amsmath,amssymb}
\usepackage{xcolor}
\usepackage{booktabs}
\usepackage{array}
\usepackage{enumitem}
\usepackage{titlesec}
\usepackage{fancyhdr}
\usepackage{tcolorbox}
\usepackage[hidelinks]{hyperref}

% ---------- colours ----------
\definecolor{accent}{RGB}{0,92,153}
\definecolor{rowgray}{gray}{0.95}

% ---------- sections ----------
\titleformat{\section}{\large\bfseries\color{accent}}{\thesection}{0.6em}{}[\vspace{-0.5ex}\titlerule]
\titlespacing*{\section}{0pt}{2.2ex}{1.2ex}
\titleformat{\subsection}{\normalsize\bfseries}{\thesubsection}{0.6em}{}
\titlespacing*{\subsection}{0pt}{1.8ex}{0.8ex}

\setlength{\parindent}{0pt}
\setlength{\parskip}{0.6ex}
\renewcommand{\arraystretch}{1.25}
\setlist[itemize]{nosep,leftmargin=1.4em}

\pagestyle{fancy}
\fancyhf{}
\fancyhead[L]{\small ARob 2026/27 - Lab 2, Sec\c{c}\~ao 5}
\fancyhead[R]{\small Pedro}
\fancyfoot[C]{\small\thepage}
\renewcommand{\headrulewidth}{0.3pt}

\newtcolorbox{caveat}{colback=rowgray,colframe=gray!60,boxrule=0.4pt,
  left=5pt,right=5pt,top=3pt,bottom=3pt}

\newcommand{\code}[1]{\texttt{#1}}
\newcommand{\Sk}{S}          % skew-symmetric matrix
\newcommand{\wm}{\omega_m}

\begin{document}

\begin{center}
{\Large\bfseries Filtro de Batista et al.\ no Lab 2 (Sec\c{c}\~ao 5)}\\[0.6ex]
{\large Apontamento sobre a l\'ogica, as escolhas e a valida\c{c}\~ao}\\[0.6ex]
{\small 2 de novembro de 2026}
\end{center}

Este apontamento explica a l\'ogica do que foi implementado, porque foi feito assim, o que foi validado e o que fica por verificar.
Os n\'umeros reportados saem todos de execu\c{c}\~oes dos scripts da pasta \code{labs-26-27/Lab2/Filter/} (comandos na sec\c{c}\~ao~\ref{sec:repro}).

%% ------------------------------------------------------------------
\section{O que foi feito}

Implementou-se o filtro de Kalman do paper de Batista, Silvestre e Oliveira (\emph{Int.\ J.\ Control}, 2011), para o sistema nominal~(2), com estima\c{c}\~ao parcial da atitude (roll e pitch) e do bias dos tr\^es rate gyros.
O filtro est\'a em Simulink, em tempo cont\'inuo, e foi testado de duas formas.
Primeiro sobre os dados reais das experi\^encias A a E (modelo de \emph{Replay} do DevKit).
Depois sobre uma simula\c{c}\~ao com verdade conhecida, onde o erro absoluto se pode medir.

%% ------------------------------------------------------------------
\section{O modelo e o filtro}

A atitude obedece a $\dot R = R\,\Sk(\omega)$, com $R$ a rota\c{c}\~ao do corpo para o referencial inercial.
O aceler\'ometro mede, em eixos do corpo, um vetor constante no referencial inercial (a gravidade): $y = R^\top c$.
Derivando, e usando $\wm = \omega + b$ (gyro com bias), obt\'em-se
\[
\dot y = -\Sk(\omega)\,y = -\Sk(\wm)\,y - \Sk(y)\,b .
\]
O estado \'e $x = [\,y;\ b\,] \in \mathbb{R}^6$, com o bias constante:
\[
\dot x = A(t)\,x,\qquad
A = \begin{bmatrix} -\Sk(\wm) & -\Sk(y) \\ 0 & 0 \end{bmatrix},\qquad
y = C x,\quad C = \begin{bmatrix} I & 0 \end{bmatrix}.
\]
O sistema \'e n\~ao linear, mas $A(t)$ s\'o depende de $y$ e de $\wm$, que s\~ao ambos medidos.
Por isso pode tratar-se como linear vari\'avel no tempo e aplicar o Kalman--Bucy cl\'assico:
\[
\dot{\hat x} = A\hat x + K\,(y - C\hat x),\qquad K = P C^\top \Theta^{-1},
\]
\[
\dot P = A P + P A^\top + \Xi - P C^\top \Theta^{-1} C P .
\]
$\Xi$ e $\Theta$ s\~ao as intensidades do ru\'ido de processo e de medida.
O ganho tem de ser recalculado a cada instante porque $A(t)$ muda (n\~ao \'e o ganho estacion\'ario da sec\c{c}\~ao~4 do guia).

\textbf{Observabilidade} (Teoremas 3.3 e 3.4 do paper): o sistema \'e observ\'avel se $y$ n\~ao for constante, e a estabilidade do erro precisa de que $y$ continue a variar (excita\c{c}\~ao persistente).
Com o drone parado, o bias na dire\c{c}\~ao da gravidade (o gyro $z$) n\~ao \'e observ\'avel.

%% ------------------------------------------------------------------
\section{Implementa\c{c}\~ao}

Ficheiros em \code{labs-26-27/Lab2/Filter/}:
\begin{itemize}
  \item \code{ARDroneReplay\_BatistaKF.slx}: o filtro sobre os dados gravados.
  \item \code{ARDroneSim\_BatistaKF.slx}: o mesmo filtro sobre uma experi\^encia sint\'etica com verdade conhecida.
  \item \code{run\_batista\_kf.m} e \code{run\_batista\_kf\_sim.m}: correm cada modelo, abrem os scopes e imprimem um resumo.
  \item \code{build\_batista\_kf\_model.m}: regenera os dois modelos (o filtro \'e um bloco partilhado).
\end{itemize}

Escolhas de implementa\c{c}\~ao:
\begin{itemize}
  \item \textbf{Tempo cont\'inuo.} Segue o paper. $\hat x$ e $P$ s\~ao integrados por blocos \emph{Integrator}, com solver de passo fixo \code{ode1} a 5~ms (o per\'iodo dos dados e o solver do modelo original do DevKit).
  \item \textbf{Unidades.} O aceler\'ometro vem em mg e passa a m/s$^2$ ($\times 9{,}81/1000$). Os gyros v\^em em deg/s e passam a rad/s. Os \^angulos mostram-se em graus.
  \item \textbf{Sinais.} Neste sensor $a_z = -g$ em repouso. Por isso $\phi = \operatorname{atan2}(-a_y,-a_z)$ e $\theta = \arcsin(a_x/\lVert y\rVert)$. A f\'ormula $\operatorname{atan2}(a_y,a_z)$ d\'a $\phi + 180^\circ$. O $\wm$ entra diretamente dos gyros: $g_x$ e $g_y$ t\^em correla\c{c}\~ao de $+0{,}97$ a $+0{,}98$ e $+0{,}94$ a $+0{,}98$ com as derivadas de $\phi$ e $\theta$.
  \item \textbf{\^Angulos.} Calculam-se a partir do vetor filtrado $\hat x_1$, e n\~ao do aceler\'ometro cru.
  \item \textbf{Condi\c{c}\~oes iniciais.} $x_0 = [0,0,-g,0,0,0]$ (nivelado, bias nulo) e $P_0 = \operatorname{blkdiag}(I,\,0{,}01\,I)$. N\~ao se usa a primeira amostra, porque o primeiro pacote do \code{navdata} \'e todo zeros.
\end{itemize}

%% ------------------------------------------------------------------
\section{Afina\c{c}\~ao}

\[
\Xi = \operatorname{blkdiag}\!\big((g\cdot 0{,}035)^2\,T_s\,I_3,\ 10^{-5} I_3\big),\qquad
\Theta = 10^{3}\cdot 0{,}07^2\,T_s\,I_3,\qquad T_s = 5~\text{ms}.
\]
\begin{itemize}
  \item \textbf{Ru\'ido de medida $\Theta$.} $0{,}07$~m/s$^2$ \'e o desvio do aceler\'ometro em hover (experi\^encia C: 0,059 a 0,077). Multiplica-se por $T_s$ para passar de vari\^ancia amostrada a intensidade cont\'inua, como no \code{kalman\_filters.m} da sec\c{c}\~ao~4.
  \item \textbf{Fator 1000 em $\Theta$.} Em voo o erro do aceler\'ometro n\~ao \'e ru\'ido branco, \'e vibra\c{c}\~ao e acelera\c{c}\~ao linear. Por isso a medida merece menos confian\c{c}a do que o ru\'ido sugere.
  \item \textbf{Ru\'ido de processo $\Xi$.} O ru\'ido dos gyros entra em $\dot y$ multiplicado por $y$, que tem m\'odulo $\approx g$, da\'i $(g\sigma_g)^2 T_s$ com $\sigma_g = 0{,}035$~rad/s (hover). O $10^{-5}$ \'e a intensidade do passeio aleat\'orio do bias.
  \item \textbf{Como foram escolhidos.} O fator 1000 e o $10^{-5}$ v\^em de um varrimento (fator entre 1 e $10^5$, intensidade do bias entre $10^{-8}$ e $10^{-4}$) nas experi\^encias D e E, a minimizar o erro contra o estimador interno do drone.
\end{itemize}

\begin{caveat}
Esta afina\c{c}\~ao \'e heur\'istica e tem duas fraquezas assumidas.
Foi feita nos mesmos dados em que depois se avalia.
E a refer\^encia (o estimador interno do drone) n\~ao \'e a verdade.
\end{caveat}

%% ------------------------------------------------------------------
\section{Valida\c{c}\~ao}

\subsection{Simula\c{c}\~ao com verdade conhecida}

Roll de $8^\circ$ a 0,30~Hz e pitch de $6^\circ$ a 0,45~Hz (fase de $60^\circ$), yaw parado.
Frequ\^encias diferentes mant\^em $y$ a mudar de dire\c{c}\~ao (condi\c{c}\~ao de excita\c{c}\~ao do Teorema 3.4).
Gyros e aceler\'ometro saem da cinem\'atica exata (verificado numericamente: $\dot y + \Sk(\omega)y \approx 10^{-8}$, contra uma derivada de $\approx 2{,}6$).
Ru\'ido gaussiano de 0,07~m/s$^2$ no aceler\'ometro e 0,01~rad/s nos gyros, bias dos gyros de $[2,-3,1]$~deg/s (o valor da simula\c{c}\~ao do paper).
Erro depois de 10~s de transit\'orio, contra a verdade:

\begin{center}
\small
\begin{tabular}{@{}lccc@{}}
\toprule
Caso & RMS $\phi/\theta$ [$^\circ$] & Erro m\'edio $\phi/\theta$ [$^\circ$] & Bias estimado [deg/s] \\
\midrule
Por defeito & 0,05 / 0,05 & $+0{,}01$ / 0,00 & [2,01, $-3{,}01$, 0,96] \\
Bias de aceler\'ometro $[-0{,}033,\,-0{,}166,\,0]$ m/s$^2$ & 0,98 / 0,20 & $+0{,}98$ / $-0{,}19$ & [2,01, $-3{,}01$, 0,89] \\
Ru\'ido $\times 0{,}01$ & 0,03 / 0,03 & 0,00 / 0,00 & [2,00, $-3{,}00$, 0,93] \\
\bottomrule
\end{tabular}
\end{center}

O bias verdadeiro dos gyros \'e $[2,-3,1]$ deg/s.
O filtro est\'a correto: o erro \'e $\approx 0{,}05^\circ$ e o bias \'e recuperado nos tr\^es eixos.
Um bias no aceler\'ometro produz um offset constante em $\phi$ e $\theta$ que o filtro n\~ao consegue corrigir.

\subsection{Dados reais contra o estimador interno do drone}

Nos dados reais n\~ao h\'a verdade.
A compara\c{c}\~ao faz-se contra os \code{states} do drone (portas 1 e 2 do bloco de Replay), em $t > 10$~s e $t < t_{\text{fim}} - 3$~s (sem descolagem nem aterragem).
O RMS calcula-se depois de retirar o offset constante, que se reporta \`a parte:

\begin{center}
\small
\begin{tabular}{@{}lccc@{}}
\toprule
Exp. & RMS $\phi/\theta$ [$^\circ$] & Offset $\phi/\theta$ sem calibrar [$^\circ$] & Offset $\phi/\theta$ calibrado [$^\circ$] \\
\midrule
C (hover)           & 0,11 / 0,11 & $+2{,}15$ / $+1{,}03$ & $+0{,}40$ / $+0{,}57$ \\
D (passos de pitch) & 0,24 / 0,55 & $+1{,}32$ / $+0{,}73$ & $-0{,}24$ / $+0{,}71$ \\
E (miss\~ao)        & 0,40 / 0,37 & $+2{,}25$ / $+0{,}30$ & $-0{,}31$ / $+0{,}28$ \\
\bottomrule
\end{tabular}
\end{center}

Para compara\c{c}\~ao, o inclin\'ometro cru (sem filtro passa-baixo) em D d\'a 0,89 / 1,61$^\circ$.
Esta linha de base \'e favor\'avel ao filtro, porque o inclin\'ometro n\~ao foi filtrado.

\subsection{Bias artificial nos dados reais}

O bias real dos gyros \'e $\approx 0$ (m\'edia em repouso de $[-0{,}03,\,-0{,}02,\,0{,}01]$ deg/s), por isso n\~ao h\'a nada a convergir.
Como sugere o guia (4.4), injetou-se um bias artificial de $[3,-2,1]$ deg/s nos gyros da experi\^encia D.
O bias estimado no fim foi $[3{,}10,\,-2{,}01,\,2{,}85]$ deg/s: converge bem em $x$ e $y$, e pior em $z$.

%% ------------------------------------------------------------------
\section{O desvio constante de \texorpdfstring{$\phi$}{phi}}

Nos dados reais o filtro fica 1 a $2^\circ$ acima do drone em $\phi$.
A explica\c{c}\~ao \'e o bias do aceler\'ometro:
\begin{itemize}
  \item A simula\c{c}\~ao mostra o mecanismo: um bias de aceler\'ometro igual ao medido em repouso (A) d\'a exatamente o offset esperado ($+0{,}98^\circ$ em $\phi$).
  \item O filtro n\~ao o pode estimar: com uma \'unica medida vetorial, o bias do aceler\'ometro \'e indistingu\'ivel de uma inclina\c{c}\~ao.
  \item O bias varia de uma experi\^encia para outra. No ch\~ao, no in\'icio de C, D e E, o $a_y$ \'e $-0{,}29$, $-0{,}26$ e $-0{,}45$ m/s$^2$ (inclina\c{c}\~ao aparente de 1,8, 1,6 e $2{,}6^\circ$), enquanto o drone marca $\approx 0^\circ$. Em repouso na experi\^encia A era $-0{,}17$ m/s$^2$ ($1{,}0^\circ$). Calibrar com A removia s\'o cerca de $1^\circ$.
\end{itemize}

Por isso implementou-se uma calibra\c{c}\~ao opcional (\code{calibrate}, desligada por defeito): m\'edia causal de $a_x$ e $a_y$ durante os primeiros 1,5~s de cada ficheiro (o drone est\'a no ch\~ao), congelada depois.
A descolagem s\'o ocorre a partir de 2,2~s nas tr\^es experi\^encias.
S\'o se corrigem $a_x$ e $a_y$, porque $a_z$ entra na inclina\c{c}\~ao s\'o em segunda ordem.

\begin{caveat}
\textbf{Pressupostos e cuidados.}
\begin{itemize}
  \item Assume que o drone come\c{c}a nivelado. Se n\~ao estiver, a calibra\c{c}\~ao desloca a inclina\c{c}\~ao.
  \item O offset que desaparece \'e contra o zero do drone (o seu trim), que n\~ao \'e a verdade. A calibra\c{c}\~ao melhora o alinhamento com o estimador interno, mas n\~ao prova que o \^angulo absoluto esteja certo.
  \item N\~ao usa a refer\^encia: s\'o o aceler\'ometro dos primeiros 1,5~s. O \'unico par\^ametro \'e a dura\c{c}\~ao, e n\~ao foi ajustado ao resultado.
  \item N\~ao altera o RMS, nem o erro nas manobras fortes (a\'i o aceler\'ometro l\^e acelera\c{c}\~ao linear como inclina\c{c}\~ao).
  \item Recomenda-se apresentar os resultados com e sem calibra\c{c}\~ao.
\end{itemize}
\end{caveat}

%% ------------------------------------------------------------------
\section{Limita\c{c}\~oes e o que n\~ao foi verificado}

\begin{itemize}
  \item O estimador interno do drone n\~ao \'e verdade. S\'o a simula\c{c}\~ao mede erro absoluto, e usa sensores simulados.
  \item A afina\c{c}\~ao foi feita nos dados de avalia\c{c}\~ao (sec\c{c}\~ao~4).
  \item Em $\theta$ sobram 0,3 a $0{,}7^\circ$ de offset em C e D que a calibra\c{c}\~ao n\~ao explica. As causas poss\'iveis (trim do drone, efeitos em voo) n\~ao foram testadas.
  \item O bias do gyro $z$ s\'o \'e fi\'avel se o drone rodar. Em E estima $-1{,}9$ deg/s, quando o bias real \'e $\approx 0{,}2$ deg/s.
  \item \code{ode1} foi escolhido por ser o solver do modelo original e porque o Replay descodifica bytes (passos interm\'edios interpol\'a-los-iam). N\~ao se comparou com RK4.
  \item O aceler\'ometro mede acelera\c{c}\~ao linear al\'em da gravidade, o que domina o erro nas manobras fortes.
\end{itemize}

%% ------------------------------------------------------------------
\section{Como reproduzir}\label{sec:repro}

No MATLAB, em \code{Filter/}:

\begin{verbatim}
out = run_batista_kf('E');                              % dados reais, 2 scopes
out = run_batista_kf('E', calibrate=1.5);               % com calibracao do acelerometro
out = run_batista_kf('D', bias=[3 -2 1]);               % bias artificial nos gyros
out = run_batista_kf_sim;                               % verdade conhecida, 3 scopes
out = run_batista_kf_sim(accbias=[-0.033 -0.166 0]);    % com bias de acelerometro
\end{verbatim}

Op\c{c}\~oes comuns: \code{pacing=0} (o mais depressa poss\'ivel) e \code{scopes=false} (sem janelas).
Cada execu\c{c}\~ao imprime o RMS, o offset e o bias estimado.

\end{document}
